\section{Related Work}
\label{sec:related}

\paragraph{Storage-backed inference and memory hierarchies.}
FlexGen~\cite{sheng2023flexgen} distributes tensors across GPU, CPU, and disk
and optimizes throughput-oriented schedules with compression and
attention-cache management. DeepSpeed-Inference~\cite{aminabadi2022deepspeed}
also provides heterogeneous inference using CPU and NVMe capacity;
ZeRO-Inference~\cite{deepspeed2022zeroinference} describes weight hosting,
layer streaming, prefetch, and release in an official technical web page.
These are direct precedents for inference beyond device residency.
ZeRO-Infinity~\cite{rajbhandari2021zeroinfinity} establishes a related
GPU/CPU/NVMe hierarchy for training, whose optimizer-state and training
results do not directly characterize inference. At the storage-access level,
LLM in a flash~\cite{alizadeh2024llmflash} uses windowing and row-column
bundling to reduce flash weight reads and enlarge useful transfers when a
model exceeds DRAM. These works motivate managing both placement and access
patterns; \system{} focuses on the contract connecting each materialization
stage to source identity, resource claims, completion evidence, and reuse.

\paragraph{Heterogeneous execution and representation.}
PowerInfer~\cite{song2023powerinfer} exploits activation locality by executing
hot neurons on a GPU and cold neurons on a CPU. This differs from
PowerInfer-2~\cite{xue2024powerinfer2}, whose cited smartphone version uses
neuron clusters, segmented caching, and an I/O/computation pipeline beyond
resident capacity. CPU execution, storage traffic, and device-cache residency
are therefore separate comparison dimensions. Weight representation is
another independent axis: AWQ~\cite{lin2024awq} uses activation-aware scaling
to protect salient weights during low-bit quantization. Representation and
materialization scheduling can be considered separately, but equal bit widths
do not establish equal algorithms, quality, or execution costs. This citation
does not imply that \system{} implements or evaluates AWQ.

\paragraph{Expert caching, prediction, and execution location.}
Direct MoE offload systems already exploit selective expert demand.
MoE-Infinity~\cite{xue2024moeinfinity} uses activation traces to guide GPU
expert-cache replacement and prefetch for host-memory-backed inference on
personal machines. Eliseev and Mazur~\cite{eliseev2023fastmoeoffloading}
combine an LRU expert cache with speculative loading from earlier hidden
states and mixed quantization. Their speculative prefetch does not itself
change routing: a prediction miss still requires the demanded expert;
quantization has a separate numerical-quality boundary.
Fiddler~\cite{kamahori2025fiddler} instead coordinates moving expert weights
to the GPU with moving activations to CPU experts. Its execution-location
choice is not equivalent to a cache policy that always executes demanded
experts on the GPU. These mechanisms concern different combinations of cache
reuse, transfer avoidance, and heterogeneous computation, not one common
storage or execution path.

\paragraph{Routing and precision changes for expert offload.}
Pre-gated MoE~\cite{hwang2024pregatedmoe} changes gating so that one block
selects experts for the next, enabling migration to overlap current expert
computation. Its pre-gate is trained during fine-tuning while retaining
pretrained parameters; this is a model change rather than speculative
prefetch applied to an unmodified checkpoint.
HOBBIT~\cite{tang2024hobbit} dynamically loads lower-precision versions of
less-critical experts on cache misses, together with adaptive prefetch and
multidimensional caching. Its loading policy thus includes a quality--latency
tradeoff. \system{} instead studies read scheduling and upload grouping with
fixed selected-expert order and arithmetic. That distinction does not establish
superior performance or rule out complementary caching policies.

\paragraph{Demand-loaded expert caches in serving engines.}
The inspected SGLang source includes a specialized expert-pack loader and
runtime cache~\cite{sglang2026expertpack}. Cache misses issue storage reads into
staging slots, followed by event-gated host-to-device copies; preparation of
the packed representation does not eliminate those runtime reads. The loader
and argument checks restrict this path to DeepSeek-V4 Flash and text-only
Kimi-K3, require unit tensor, data, and expert parallelism, and disable CUDA
graphs and shared-expert fusion. These are boundaries of the inspected path,
not claims about every SGLang configuration or future version. We neither ran
this path nor established its release inclusion, numerical qualification, or
relative performance. Its source demonstrates that storage-backed expert
loading is not unique to \system{}; equivalent failure-custody guarantees
have not been exhaustively audited in either direction.

\paragraph{KV allocation, selection, and prefix reuse.}
KV state is a different object from checkpoint weights.
vLLM uses PagedAttention~\cite{kwon2023pagedattention} for block-based KV
allocation and sharing, and its paper includes CPU KV swapping for preempted
requests. The separately inspected CPU-weight-offload paths use host-resident
weights, including unified virtual addressing; these should not be equated
with disk-backed weight materialization. SGLang uses a radix tree for KV
prefix reuse across calls~\cite{zheng2024sglang}; that paper is distinct from
the expert-pack source snapshot above.
Pensieve~\cite{yu2025pensieve} retains conversation state across requests in
a GPU/CPU cache to avoid repeated history processing.
InfiniGen~\cite{lee2024infinigen} instead predicts important tokens and
prefetches selected KV from host memory. Importance-based approximate
selection must be distinguished from retaining complete state: its quality
conditions are not those of exact capacity management.

\paragraph{KV storage and transport.}
Mooncake~\cite{qin2025mooncake} builds a cluster-wide KV-centric serving
architecture using CPU, DRAM, SSD, and network resources, establishing an
explicit storage-backed KV precedent. CacheGen~\cite{liu2024cachegen}
compresses and streams KV with bandwidth-adaptive compression, addressing
network loading and its quality tradeoff rather than only local placement.
D{\'e}j{\`a}Vu~\cite{strati2024dejavu} provides KV streaming for prompt/token
separation, microbatch swapping, and state replication for recovery.
These systems show why reuse scope, storage tier, transfer path, and numerical
semantics must be stated separately. KV transport and replication can support
a serving architecture without being the same mechanism as checkpoint-weight
materialization or establishing its resource-lifetime guarantees.

\paragraph{Request and phase scheduling.}
Orca~\cite{yu2022orca} schedules at generation-iteration boundaries with
selective batching. Sarathi-Serve~\cite{agrawal2024sarathiserve} combines
chunked prefill and stall-free scheduling to manage the throughput--latency
tradeoff. Splitwise~\cite{patel2024splitwise} places prompt processing and
token generation on separate machines according to their resource demands;
DistServe~\cite{zhong2024distserve} separates prefill and decoding to optimize
goodput under latency constraints. At a broader placement level,
AlpaServe~\cite{li2023alpaserve} combines model parallelism and statistical
multiplexing for bursty workloads. These works choose when and where requests
or phases execute. \system{} addresses resource admission and completion
within materialization, not a new request-batching or cluster-placement
algorithm; this positioning implies no cross-system performance ordering.

\paragraph{Execution control.}
Salus~\cite{yu2020salus} separates GPU-sharing primitives, including fast
switching and memory sharing, from scheduling policy.
Clockwork~\cite{gujarati2020clockwork} builds predictable serving around
explicit control of execution and resource behavior.
Tutel~\cite{hwang2023tutel} adapts parallelism and pipelined execution to
dynamic MoE workloads. These works concern execution control rather than the
source-bound materialization lifecycle studied here. Their results do not
establish \system{} performance, predictability, or completion guarantees.

\paragraph{Bounded prefetch and staged admission.}
Informed prefetching~\cite{patterson1995informed} uses future-access hints to
coordinate prefetching with cache-buffer allocation. This precedent treats
lookahead as a consumer of finite buffers, not free early I/O.
SEDA~\cite{welsh2001seda} connects event-driven stages with explicit queues
and stage-local resource control. \system{} applies bounded admission to
materialization claims; it does not implement those systems' buffer-allocation
algorithms or adaptive controllers.

\paragraph{Publication and safe reclamation.}
Linearizability~\cite{herlihy1990linearizability} is a correctness condition
for concurrent-object histories, not a resource-reclamation mechanism or a
persistent transaction protocol. Hazard pointers~\cite{michael2004hazard}
distinguish logical removal from safe reclamation after concurrent use ends.
\system{} uses a local publication and completion protocol with owner-held
evidence through retirement. This is neither a linearizability proof nor a
persistent ACID protocol, and the reclamation analogy does not imply a
hazard-pointer implementation or a general device-reclamation theorem.

\paragraph{Measurement and comparison boundaries.}
An asynchronous interface, an independent progress worker, and hardware
overlap are distinct properties. The earlier read-scheduling comparison holds
worker, representation, and resource limits fixed, then changes when one
next-expert read is submitted. The transfer-policy comparison instead holds
the read schedule fixed and changes upload grouping. Client timing measures
the resulting behavior; correlated file-call and kernel intervals separately
measure observed read--compute overlap. Selected-process API counts quantify
observed synchronization calls, not a complete causal decomposition. Neither
a worker, a broad host-range intersection, nor fewer synchronization calls
alone establishes device overlap or faster serving. A fair external comparison
must also match representation, arithmetic, physical memory pools, workload,
and quality criterion. The unmatched resident-system observations preserved
with the artifacts do not satisfy those conditions and support no comparative
speedup or ranking here.

\begin{table}[!htbp]
\centering
\small
\begin{tabularx}{\linewidth}{@{}>{\raggedright\arraybackslash}p{0.20\linewidth}>{\raggedright\arraybackslash}X>{\raggedright\arraybackslash}X@{}}
\toprule
System & Resource focus & Principal mechanism \\
\midrule
FlexGen & GPU, CPU, disk & Tensor placement, batching, compression \\
ZeRO-Inference & CPU or NVMe to GPU & Layer streaming and prefetch \\
PowerInfer & GPU--CPU residency & Hot/cold neuron execution \\
PowerInfer-2 & Smartphone memory/storage & Neuron clusters and I/O pipeline \\
SGLang expert pack & Storage to GPU expert cache & Demand reads, staging, event-gated copies \\
vLLM & Serving KV residency; host-weight offload & PagedAttention; separate resident offload paths \\
\bottomrule
\end{tabularx}
\caption{Mechanism-level positioning from the cited papers and inspected source.
The SGLang entry concerns its specialized expert-pack path, not all adapters.
These are not matched performance measurements.}
\label{tab:prior-work}
\end{table}

\system{} contributes a profile-scoped account of source binding, stage-local
admission, publication, and retirement, together with measurements that keep
logical and physical observations separate. We have not exhaustively audited
the custody or failure semantics of the systems above and do not claim that
they lack equivalent mechanisms. Model-registry entries and architectural
capability are not substitutes for qualified execution: in particular,
\system{} has no tested DeepSeek-V4 serving result. The empirical contribution
is controlled internal comparisons and their explicit resource and completion
scope, not a first-system or state-of-the-art claim.
